\documentclass[10pt,a4paper]{amsart}
\usepackage[margin=19mm]{geometry}
\usepackage{amsmath,amssymb,mathtools}
\usepackage[scr=rsfs,cal=euler]{mathalfa}
\usepackage{xfrac}
\usepackage[unicode,hidelinks,bookmarks=false]{hyperref}
\newcommand{\ie}{i.e.\ }
\newcommand{\cf}{cf.\ }
\newcommand{\resp}{respectively\ }
\newcommand{\Subsection}{\subsection}
\providecommand{\nobreakdash}{\nobreak\hbox{-}}
\setlength{\emergencystretch}{2em}
\multlinegap0pt
\usepackage{amssymb}
\usepackage{mathtools}
\usepackage[shortlabels]{enumitem}
\usepackage{xcolor}
\usepackage{tikz}
\newtheorem{theorem}{Theorem}
\newtheorem{remark}{Remark}
\newtheorem{lemma}{Lemma}
\newcommand{\K}{\mathcal{K}}
\newcommand{\M}{\mathcal{M}}
\newcommand{\N}{\mathbb{N}}
\def\Q{{\mathbb Q}}
\newcommand{\R}{\mathbb{R}}
\def\ge{\geqslant}
\def\le{\leqslant}
\title{Equidistribution and thermodynamics at infinity}
\author{Godofredo Iommi, Felipe Riquelme, An\'ibal Velozo}
\date{Source: arXiv:2608.19457v1 (CC BY 4.0)}
\begin{document}
\maketitle

\noindent\emph{Source and licence.} Equidistribution and thermodynamics at infinity. Version arXiv:2608.19457v1 (CC BY 4.0), distributed by arXiv under the Creative Commons Attribution 4.0 International licence, http://creativecommons.org/licenses/by/4.0/, as recorded in the arXiv licence metadata for this exact version and shown on the abstract page https://arxiv.org/abs/2608.19457v1. Full licence notice: \texttt{licenses/arxiv-2608.19457-CC-BY-4.0.txt}.\par\smallskip

\noindent\textit{Editorial scope.} Counted unit: the article's theorem thm:prob (source0.tex lines 505--521). The article's complete original proof of it is delivered with it (source0.tex lines 2405--2655, one \texttt{proof} environment of 251 lines, which the article places away from the statement). No mathematics is written, repaired or re-derived: the statement and the proof are the article's own lines, in the article's own notation, and no retained line is substituted or re-flowed.  Retained uncounted as the source-local context the proof needs: lines 383--388; lines 1007--1008; lines 1408--1409.  Nothing was omitted from any retained span, so the package carries no omission record.  No retained line contains a citation, so the wrapper carries no bibliography and no bibliography entry.  No retained line contains a figure environment, an image-inclusion command or drawing code, so no image is delivered and the package declares no asset dependency.  Editorial formatting only, in addition: an AMS \texttt{amsart} preamble that restates the article's own declarations and macros used by the retained lines, namely the environments \texttt{theorem}, \texttt{remark}, \texttt{lemma} on separate counters, together with the standard packages those lines need.  The retained environments print in this document as Theorem 1, Remark 1, Lemma 1 in order of appearance, whereas the article prints the same items with the numbering of its own full document.  The complete native source of the article as distributed in the arXiv e-print for this exact version is bundled unchanged as \texttt{sources/208090/source0.tex}.
\begin{theorem}\label{thm_ldp2} Let $(\Sigma,\sigma)$ be a topologically mixing countable Markov shift and  $\tau:\Sigma\to\R$ a function of summable variations satisfying $\inf\tau>0$. Denote by  $(Y,\Theta)$ the associated suspension semi-flow and assume that it has finite entropy. Let $f\in \mathrm{C}_b(Y)$ be such that  $\Delta_f$ has summable variations. Let $a\in \N$, $d>0$ and $\mathcal{K}$ be a closed subset of $\M_{\le 1}(\Theta)$. Then,
\begin{align*}
\limsup_{T \to \infty} \frac{1}{T} \log  &\left(\frac{\sum_{(x;s)\in \mathrm{FPer}_a(T-d,T), \nu_{x}\in \mathcal{K}}\exp\big(\int_0^{s}f(\theta_t(x,0))dt\big)}{\sum_{(x;s)\in \mathrm{FPer}_a(T-d,T)}\exp\big(\int_0^{s}f(\theta_t(x,0))dt\big)}\right) \leq - \inf \left\{I_f(\nu) : \nu \in \K \right\}.
\end{align*}
If in addition $P_{\Theta}(f)\ge0$ then we can replace $\mathrm{FPer}_a(T-d,T)$ with $\mathrm{FPer}_a(T)$.  
\end{theorem}

\begin{remark}\label{rem:convprob} Let $(\nu_n)_n$ and $\nu$ be measures in  $\M(\Theta)$. Set $\mu_n=(AK)^{-1}(\nu_n)$ and $\mu=(AK)^{-1}(\nu)$. Note that $(\nu_n)_n$ converges on cylinders to $\nu$ if and only if $(\mu_n)_n$ converges on cylinders to $\mu$, and $\lim_{n\to\infty}\int \tau d\mu_n=\int \tau d\mu$. Since the cylinder topology coincides with the weak$^\ast$ topology in the absence of escape of mass, this statement can equivalently be formulated in terms of weak$^\ast$ convergence.
\end{remark}

\begin{lemma} \label{lem:beta} If $\nu\in \M_{\le1}(\Theta)$ satisfies that $I_f(\nu)=0$, then $\nu=\nu_f$. 
\end{lemma}

\begin{theorem}\label{thm:prob}
For every sufficiently small $\epsilon>0$, there exists $r_\epsilon>0$ such that, for all sufficiently large $q\in\mathbb N$,
\[
\frac{1}{\Phi(q)}
\#
\left\{
\frac{p}{r}\in\Q(q):
\left|
\frac{\ell(p/r)}{2\log r}
-
\frac{6\log 2}{\pi^2}
\right|
\geq \epsilon
\right\}
\leq q^{-2r_\epsilon}.
\]
\end{theorem}

\begin{proof}[Proof of Theorem~\ref{thm:prob}]
We use the symbolic coding of the Gauss map. Let
\(\Sigma=\mathbb N^{\mathbb N}\) be the full shift and let
\(\pi:\Sigma\to(0,1)\) be the continued-fraction coding map. Define $\widehat\tau(x):=\log |G'(\pi x)|$ and
$\tau:=\frac12\left(\widehat\tau+\widehat\tau\circ\sigma\right)$. Then $\tau$ has summable variations and is bounded away from zero. Moreover, for every periodic point $x$ of period $m$, we have $S_m\tau(x)=S_m\widehat\tau(x)$.
Let $Y$ be the suspension over $(\Sigma,\sigma)$ with roof function $\tau$, and let $\Theta=(\Theta_t)_{t\ge0}$ be the corresponding suspension semi-flow. Recall that this suspension flow has topological entropy equal to $1$, and its measure of maximal entropy is $\nu_{mme}=AK(m_G),$
where $m_G$ is the Gauss measure in symbolic coordinates. Furthermore,
\[
\int \tau\,dm_G
=
\int \widehat\tau\,dm_G
=
\int \log |G'|\,d\mu_G
=
\frac{\pi^2}{6\log 2}.
\]
Set $\gamma:=\frac{6\log 2}{\pi^2}$. For a finite word $w=(a_1,\ldots,a_m)$, let $x_w=\overline{(a_1,\ldots,a_m)}$
be the corresponding periodic point, and $s(w):=S_m\tau(x_w).$
Let
\[
\mu_w:=\frac1m\sum_{j=0}^{m-1}\delta_{\sigma^j x_w}
\quad \text{ and } \quad \nu_w:=AK(\mu_w)
\]
be the corresponding invariant probability measure for the suspension flow. Then
\[
\int\tau\,d\mu_w=\frac{s(w)}{m},
\quad \text{ and therefore } \quad
\frac{m}{s(w)}
=
\left(\int\tau\,d\mu_w\right)^{-1}.
\]
Fix $\epsilon>0$. Define
\[
\mathcal E_\epsilon
:=
\left\{
\nu_w:
\left|
\frac{|w|}{s(w)}-\gamma
\right|
\ge \frac{\epsilon}{2}
\right\},
\]
and let
\[
\mathcal K_\epsilon
:=
\overline{\mathcal E_\epsilon}
\subset \mathcal M_{\le1}(\Theta),
\]
where the closure is taken in the cylinder topology. Thus
\(\mathcal K_\epsilon\) is closed by definition.
We claim that $\nu_{mme}\notin \mathcal K_\epsilon$.
Indeed, suppose by contradiction that there exists a sequence $(w_n)_n$ such that $\nu_{w_n}\to\nu_{mme}$
in the cylinder topology and for every $n \in \N$ we have
\[
\left|
\frac{|w_n|}{s(w_n)}-\gamma
\right|
\ge \frac{\epsilon}{2}
\]
Since both $\nu_{w_n}$ and $\nu_{mme}$ are probability measures, Remark~\ref{rem:convprob} implies that
$\mu_{w_n}\to m_G$ on cylinders and
\[
\int\tau\,d\mu_{w_n}
\longrightarrow
\int\tau\,dm_G .
\]
Consequently,
\[
\frac{|w_n|}{s(w_n)}
=
\left(\int\tau\,d\mu_{w_n}\right)^{-1}
\longrightarrow
\left(\int\tau\,dm_G\right)^{-1}
=
\gamma,
\]
which contradicts the definition of \(\mathcal E_\epsilon\). Hence
\(\nu_{mme}\notin\mathcal K_\epsilon\).

By Lemma~\ref{lem:beta}, or equivalently by the positivity of the rate
function away from the unique measure of maximal entropy, we have
\[
\beta_\epsilon:=\inf_{\nu\in\mathcal K_\epsilon} I_0(\nu)>0.
\]
Applying Theorem~\ref{thm_ldp2} with potential \(0\), and using the
full-shift/BIP version of the periodic-orbit estimates, there exist constants
\(C_\epsilon>0\), \(\rho_\epsilon>0\), and \(T_\epsilon>0\) such that, for all
\(T\ge T_\epsilon\),
\begin{equation}\label{eq:ldp-cumulative-prob}
\#\left\{
(x;s)\in\mathrm{FPer}(T):
\nu_x\in\mathcal K_\epsilon
\right\}
\le
C_\epsilon e^{(1-\rho_\epsilon)T}.
\end{equation}
Here \(\mathrm{FPer}(T)\) denotes periodic data of suspension length at most
\(T\). This cumulative estimate follows from the window estimate in
Theorem~\ref{thm_ldp2} by summing over intervals of fixed size; the resulting
polynomial factor is absorbed after decreasing \(\rho_\epsilon\), if
necessary.

We now compare suspension length with denominators of rationals. Let
\(p/r\in(0,1)\), \((p,r)=1\), and choose its canonical finite continued
fraction expansion
\[
\frac pr=[a_1,\ldots,a_m],
\]
so that \(m=\ell(p/r)\). Let \(w=(a_1,\ldots,a_m)\). By bounded distortion for
continued-fraction cylinders, there exists \(C_0>0\), independent of \(p/r\),
such that
\begin{equation}\label{eq:susp-denom-comparison}
\left|s(w)-2\log r\right|\le C_0 .
\end{equation}
Let $\tau_0:=\inf\tau>0.$ Since $s(w)\ge m\tau_0$, we obtain
\[
\begin{aligned}
\left|
\frac{m}{s(w)}
-
\frac{m}{2\log r}
\right|
&=
m
\left|
\frac{2\log r-s(w)}{s(w)\,2\log r}
\right|  
&\le
\frac{C_0m}{s(w)\,2\log r}
\le
\frac{C_0}{2\tau_0\log r}.
\end{aligned}
\]
Therefore, there exists \(R_\epsilon\ge1\) such that, whenever
\(r\ge R_\epsilon\),
\[
\left|
\frac{m}{2\log r}-\gamma
\right|
\ge \epsilon
\quad\Longrightarrow\quad
\left|
\frac{m}{s(w)}-\gamma
\right|
\ge \frac{\epsilon}{2}.
\]
Equivalently, for \(r\ge R_\epsilon\),
\[
\left|
\frac{\ell(p/r)}{2\log r}
-
\frac{6\log 2}{\pi^2}
\right|
\ge \epsilon
\quad\Longrightarrow\quad
\nu_w\in\mathcal K_\epsilon .
\]
Now let
\[
\Q(q)
=
\left\{
\frac pr\in(0,1):
1\le p<r\le q,\ (p,r)=1
\right\}
\]
and
\[
\Phi(q)=\#\Q(q)=\sum_{r\le q}\varphi(r).
\]
If \(p/r\in\Q(q)\), then \(r\le q\), and hence
\[
s(w)\le 2\log r+C_0\le 2\log q+C_0.
\]
Thus the set
\[
\left\{
\frac pr\in\Q(q):
\left|
\frac{\ell(p/r)}{2\log r}-\gamma
\right|
\ge\epsilon
\right\}
\]
is contained, up to the finitely many rationals with \(r<R_\epsilon\), in the
set of periodic data satisfying
\[
s(w)\le 2\log q+C_0
\qquad\text{and}\qquad
\nu_w\in\mathcal K_\epsilon .
\]
The finitely many exceptional rationals contribute \(O_\epsilon(1)\). Hence,
by \eqref{eq:ldp-cumulative-prob}, for all sufficiently large \(q\),
\[
\begin{aligned}
\#\left\{
\frac pr\in\Q(q):
\left|
\frac{\ell(p/r)}{2\log r}-\gamma
\right|
\ge\epsilon
\right\}
&\le
O_\epsilon(1)
+
C_\epsilon
\exp\left((1-\rho_\epsilon)(2\log q+C_0)\right) 
&\le
C'_\epsilon q^{2(1-\rho_\epsilon)} .
\end{aligned}
\]

Finally, since
\[
\Phi(q)=\sum_{r\le q}\varphi(r)\asymp q^2,
\]
there exists \(c>0\) such that \(\Phi(q)\ge cq^2\) for all sufficiently large
\(q\). Therefore
\[
\frac{1}{\Phi(q)}
\#\left\{
\frac pr\in\Q(q):
\left|
\frac{\ell(p/r)}{2\log r}-\gamma
\right|
\ge\epsilon
\right\}
\le
C''_\epsilon q^{-2\rho_\epsilon}.
\]
After decreasing the exponent once more, we obtain a constant
\(r_\epsilon>0\) such that, for all sufficiently large \(q\),
\[
\frac{1}{\Phi(q)}
\#
\left\{
\frac pr\in\Q(q):
\left|
\frac{\ell(p/r)}{2\log r}
-
\frac{6\log 2}{\pi^2}
\right|
\ge\epsilon
\right\}
\le
q^{-2r_\epsilon}.
\]
This proves the theorem.
\end{proof}

\end{document}
